\documentclass[10.5pt]{article}
\usepackage{amsmath, amsfonts, amssymb,amsthm}
\usepackage[includeheadfoot]{geometry} % For page dimensions
\usepackage{fancyhdr}
\usepackage{enumerate} % For custom lists
\usepackage{tikz-cd}
\usepackage{graphicx}

\fancyhf{}
\lhead{MAT1301 hw7}
\rhead{Tighe McAsey - 1008309420}
\pagestyle{fancy}

% Page dimensions
\geometry{a4paper, margin=1in}

\theoremstyle{definition}
\newtheorem{pb}{}
\usepackage{tikz-cd, stackengine}

% Commands:
\tikzset{
  curarrow/.style={
  rounded corners=8pt,
  execute at begin to={every node/.style={fill=red}},
    to path={-- ([xshift=-50pt]\tikztostart.center)
    |- (#1) node[fill=white] {$\scriptstyle \delta$}
    -| ([xshift=50pt]\tikztotarget.center)
    -- (\tikztotarget)}
    }
}

\newcommand{\set}[1]{\{#1\}}
\newcommand{\gen}[1]{\langle#1\rangle}
\newcommand{\abs}[1]{\left\vert#1\right\vert}
\newcommand{\norm}[1]{\lvert\lvert#1\rvert\rvert}
\newcommand{\tand}{\text{ and }}
\newcommand{\tor}{\text{ or }}
\newcommand{\pd}{\frac{\partial}{\partial x_j}}
\newcommand{\px}{\frac{\partial}{\partial x}}
\newcommand{\py}{\frac{\partial}{\partial y}}
\newcommand{\pz}{\frac{\partial}{\partial z}}
\newcommand{\ppx}{\frac{\partial^2}{\partial x^2}}
\newcommand{\ppy}{\frac{\partial^2}{\partial y^2}}
\newcommand{\ppz}{\frac{\partial^2}{\partial z^2}}
\newcommand{\hess}{\operatorname{Hess}}

\begin{document}
    \begin{pb}
        
    \end{pb}
    \begin{pb}
        I will prove that \(\tilde{H}_k(BM((X_j,x_j)^m_1)) \cong \tilde{H}_k(BM((X_j,x_j)_1^{m-1}))\oplus \tilde{H}_k(X_m)\), in which case by induction we get \(\tilde{H}_k(BM((X_j,x_j)^m_1)) \cong \bigoplus_{j=1}^m \tilde{H}_k(X_j)\). To prove this induction step, we use Mayer Vietoris, to specify the open sets denote the branches of the tree \(Y\) as \(y_j \in I_j\), and let \(U_0\) be a small open neighborhood of \(y_0\), i.e. \(U_0 = \bigcup_1^m[0,\epsilon)_j\)  where \([0,\epsilon_j) \subset I_j\) and \(\epsilon << 1\). for each \((X_j,x_j)\) let \(U_j = U_0 \cup I_j \cup X_j\), then clearly \(U_j \simeq X_j\) by a deformation retract. We define \(V = \bigcup_1^{m-1} U_j\), then \(V \simeq BM((X_j,x_j)_1^{m-1})\) by a deformation retract, and denote \(U = U_m\). Finally, we see that \(U \cap V \simeq \set{y_0}\) by a deformation retract, so that applying Mayer Vietoris (on reduced homology) yields the exact sequence
        \begin{equation*}
            \begin{tikzcd}[arrow style=math font,cells={nodes={text height=2ex,text depth=0.75ex}}]
                \cdots & \tilde{H}_{k-1}(U) \oplus \tilde{H}_{k-1}(V) \arrow[l] \arrow[draw=none]{d}[name=Y, shape=coordinate]{} & \arrow[l] \tilde{H}_{k-1}(U \cap V) \\
                \tilde{H}_{k}(BM) \arrow[curarrow=Y]{urr}{} & \tilde{H}_{k}(U) \oplus \tilde{H}_{k}(V) \arrow[l] \arrow[draw=none]{d}[name=Z,shape=coordinate]{} & \tilde{H}_{k}(U \cap V) \arrow[l] \\
                \tilde{H}_{k+1}(BM) \arrow[curarrow=Z]{urr}{} & \tilde{H}_{k+1}(U) \oplus \tilde{H}_{k+1}(V) \arrow[l] & \cdots \arrow[l]
            \end{tikzcd}
        \end{equation*}
        But since we are working with reduced homology, and and \(U \cap V\) is connected, we get \(\tilde{H}_k(U \cap V) = \tilde{H}_k(\set{y_0}) = 0\) for all \(k\), thus the above sequence simplifies to
         \begin{equation*}
            \begin{tikzcd}[arrow style=math font,cells={nodes={text height=2ex,text depth=0.75ex}}]
                \cdots & \tilde{H}_{k-1}(U) \oplus \tilde{H}_{k-1}(V) \arrow[l] \arrow[draw=none]{d}[name=Y, shape=coordinate]{} & \arrow[l] 0 \\
                \tilde{H}_{k}(BM) \arrow[curarrow=Y]{urr}{} & \tilde{H}_{k}(U) \oplus \tilde{H}_{k}(V) \arrow[l] \arrow[draw=none]{d}[name=Z,shape=coordinate]{} & 0 \arrow[l] \\
                \tilde{H}_{k+1}(BM) \arrow[curarrow=Z]{urr}{} & \tilde{H}_{k+1}(U) \oplus \tilde{H}_{k+1}(V) \arrow[l] & \cdots \arrow[l]
            \end{tikzcd}
        \end{equation*}
        So that by exactness we have for all \(k\), \[\tilde{H}_k(BM((X_{j},x_j)_1^m)) \cong \tilde{H}_k(U) \oplus \tilde{H}_k(V) \cong \tilde{H}_k(X_j) \oplus \tilde{H}_k(BM((X_{j},x_j)_1^{m-1}))\]
        where we use homotopy invariance of homology for the second isomorphism. \qed
    \end{pb}
    \begin{pb}
        Let \(CX = X \times I/ X\times \set{1}\) denote the mapping cone of \(X\), then it is immediate that \((X,CX)\) is a good pair (\(X\) includes along the zero coordinate) implying we get \(\tilde{H}_k(CX,X) \cong \tilde{H}_k(\Sigma X)\), and that \(CX\) is contractible to \(\set{1}\). Using this, we get that \(\tilde{H}_k(CX) = 0\) for all \(k\), and the following LES of the pair
        \begin{equation*}
            \begin{tikzcd}[arrow style=math font,cells={nodes={text height=2ex,text depth=0.75ex}}]
                \cdots & 0 \arrow[l] \arrow[draw=none]{d}[name=Y, shape=coordinate]{} & \arrow[l] \tilde{H}_{k-1}(X) \\
                \tilde{H}_{k}(\Sigma X) \arrow[curarrow=Y]{urr}{} & 0 \arrow[l] \arrow[draw=none]{d}[name=Z,shape=coordinate]{} & \tilde{H}_k(X) \arrow[l] \\
                \tilde{H}_{k+1}(\Sigma X) \arrow[curarrow=Z]{urr}{} & 0 \arrow[l] & \cdots \arrow[l]
            \end{tikzcd}
        \end{equation*}
        Exactness gives us that \(\tilde{H}_{k+1}(\Sigma X) \cong \tilde{H}_k(X)\), and \(\tilde{H}_0(\Sigma X) = 0\). \qed
    \end{pb}
    \begin{pb} \textbf{(a)}
        We use the following open sets and Mayer Vietoris

        \textcolor{red}{Inssert Image Here}

        Then \(U \simeq S^1 \vee S^1, V \simeq \set{\text{pt.}}\), and \(U \cap V \simeq S^1\), thus \(\tilde{H}_k(U) \cong \tilde{H}_k(S^1 \vee S^1),  0 = \tilde{H}_k(V)\) for all \(k\), and \(\tilde{H}_k(U \cap V) \cong \tilde{H}_k(S^1)\) for all \(k\). We first compute \(\tilde{H}_{k}(S^1 \vee S^1)\) also using Mayer Vietoris, for choosing \(U,V\) to be contractible neighborhoods of either copy of \(S^1\), so that \(U \cap V \simeq \set{\text{pt}}\) is a contractible neighborhood of the wedge point. Whence \(\tilde{H}_{k}(U \cap V) = 0\) for all \(k\), and \(\tilde{H}_{k}(U) \cong \tilde{H}_{k}(S^1) \cong \tilde{H}_{k}(V)\) for all \(k\). This gives the following long exact sequence using Mayer Vietoris:
        \begin{equation*}
            \begin{tikzcd}[arrow style=math font,cells={nodes={text height=2ex,text depth=0.75ex}}]
                \cdots & \tilde{H}_{k-1}(S^1) \oplus \tilde{H}_{k-1}(S^1) \arrow[l] \arrow[draw=none]{d}[name=Y, shape=coordinate]{} & \arrow[l] 0 \\
                \tilde{H}_{k}(S^1 \vee S^1) \arrow[curarrow=Y]{urr}{} & \tilde{H}_{k}(S^1) \oplus \tilde{H}_{k}(S^1) \arrow[l] \arrow[draw=none]{d}[name=Z,shape=coordinate]{} & 0 \arrow[l] \\
                \tilde{H}_{k+1}(S^1 \vee S^1) \arrow[curarrow=Z]{urr}{} & \tilde{H}_{k+1}(S^1) \oplus \tilde{H}_{k+1}(S^1) \arrow[l] & \cdots \arrow[l]
            \end{tikzcd}
        \end{equation*}
        So that \[\tilde{H}_k(S^1 \vee S^1) \cong \tilde{H}_{k}(S^1) \oplus \tilde{H}_{k}(S^1) \cong \begin{cases}
            \mathbb{Z} \oplus \mathbb{Z} & k = 1 \\
            0 & \text{else}
        \end{cases}\]
        The last key ingredient we need is to compute \(\iota_*: \tilde{H}_1(U \cap V) \to \tilde{H}_1(U) \oplus \tilde{H}_1(V)\). To do so, we use that \(\sigma: \Delta^1 = I \to S^1\) via \(x \mapsto e^{2\pi i x}\) has \([\sigma] = 1\). And moreover, by the above computation for \(S^1 \vee S^1\), we may identify the generators for \(\tilde{H}_1(S^1 \vee S^1)\) as this \(\sigma\) on either copy of \(S^1\). Now we may use that the map \(\sigma: I \to U \cap V\), which is the map with winding number one onto the copy of \(S^1\) we are deformation retracting onto, satisfies that its image under the retract identifying \(U\) with \(S^1 \vee S^1\) is \(\sigma_1 + \sigma_2 - \sigma_1 - \sigma_2 = 0\), as in the following picture

        \textcolor{red}{Include Image Here}

        This of course implies that \(\iota_*: \tilde{H}_1(S^1) \to \tilde{H}_1(S^1\vee S^1) \oplus 0\) is the zero map. Whence we get the following Mayer-Vietoris sequence for \(T^2\):
        \begin{equation*}
            \begin{tikzcd}[arrow style=math font,cells={nodes={text height=2ex,text depth=0.75ex}}]
                \tilde{H}_0(T^2) & \tilde{H}_0(S^1 \vee S^1) \oplus 0 \arrow[l] \arrow[draw=none]{d}[name=Y, shape=coordinate]{} & \arrow[l] \tilde{H}_0(S^1) \\
                \tilde{H}_1(T^2) \arrow[curarrow=Y]{urr}{} & \tilde{H}_{1}(S^1 \vee S^1) \oplus 0 \arrow[l] \arrow[draw=none]{d}[name=Z,shape=coordinate]{} & \tilde{H}_1(S^1) \arrow[l,"0"] \\
                \tilde{H}_{2}(T^2) \arrow[curarrow=Z]{urr}{} & \tilde{H}_{2}(S^1 \vee S^1) \oplus 0 \arrow[l] & \cdots \arrow[l]
            \end{tikzcd}
        \end{equation*}
        Cleaning things up further, we get all \(\tilde{H}_k(S^1 \vee S^1) = 0 = \tilde{H}_k(S^1)\) for \(k \geq 2\), so that \(\tilde{H}_k(T^2) \cong 0\) for \(k \geq 3\), and \(\tilde{H}_0(T^2) = 0\) by connectedness. Writing everything up to isomorphism now gives
        \begin{equation*}
            \begin{tikzcd}[arrow style=math font,cells={nodes={text height=2ex,text depth=0.75ex}}]
                0 & 0 \arrow[l] \arrow[draw=none]{d}[name=Y, shape=coordinate]{} & \arrow[l] 0 \\
                \tilde{H}_1(T^2) \arrow[curarrow=Y]{urr}{} & \mathbb{Z} \oplus \mathbb{Z} \arrow[l] \arrow[draw=none]{d}[name=Z,shape=coordinate]{} & \mathbb{Z} \arrow[l,"0"] \\
                \tilde{H}_{2}(T^2) \arrow[curarrow=Z]{urr}{} & 0 \oplus 0 \arrow[l] & \cdots \arrow[l]
            \end{tikzcd}
        \end{equation*}
        So that since the map \(\mathbb{Z} \to \mathbb{Z} \oplus \mathbb{Z}\) is zero, we find that \(\tilde{H}_2(T^2) \cong \mathbb{Z}\) and \(\tilde{H}_1(T^2) \cong \mathbb{Z} \oplus \mathbb{Z}\) by exactness. Combining everything we find that
        \begin{align*}
            \tilde{H}_k(T^2) \cong \begin{cases}
                \mathbb{Z} & k = 2 \\
                \mathbb{Z} \oplus \mathbb{Z} & k = 1 \\
                0 & \text{else}
            \end{cases}
        \end{align*} \qed

        \textbf{(b)} denote \(M \simeq S^1\) to be the mobius strip, then let \(U_0\) be a small neighborhood of \(\set{x_0} \times S^1\) in \(T^2\), which deformation retracts to \(\set{x_0} \times S^1\) in \(T^2\) and \(U = U_0 \cup M \simeq M\), by simply retracting \(U_0\) to \(\set{x_0} \times S^1\). Let \(V_0\) be a small neighborhood of \(\partial M\), which retracts to the boundary, then let \(V = V_0 \cup T^2 \simeq T^2\) by contracting \(V_0\), moreover \(U \cap V = U_0 \cup V_0 =\simeq\set{x_0} \times S^1\). This gives us the Mayer Vietoris sequence:
         \begin{equation*}
            \begin{tikzcd}[arrow style=math font,cells={nodes={text height=2ex,text depth=0.75ex}}]
                \cdots & \tilde{H}_{k-1}(S^1) \oplus \tilde{H}_{k-1}(T^2) \arrow[l] \arrow[draw=none]{d}[name=Y, shape=coordinate]{} & \arrow[l] \tilde{H}_{k-1}(S^1) \\
                \tilde{H}_{k}(X) \arrow[curarrow=Y]{urr}{} & \tilde{H}_{k}(S^1) \oplus \tilde{H}_{k}(T^2) \arrow[l] \arrow[draw=none]{d}[name=Z,shape=coordinate]{} & \tilde{H}_{k}(S^1) \arrow[l] \\
                \tilde{H}_{k+1}(X) \arrow[curarrow=Z]{urr}{} & \tilde{H}_{k+1}(S^1) \oplus \tilde{H}_{k+1}(T^2) \arrow[l] & \cdots \arrow[l]
            \end{tikzcd}
        \end{equation*}
        since \(\tilde{H}_k(S^1) = 0\), for \(k \geq 2\), and \(\tilde{H}_k(T^2) = 0\) for \(k \geq 3\), we find that \(\tilde{H}_k(X) = 0\) for \(k \geq 3\). Morover, \(X\) is connected, so that \(\tilde{H}_0(X) = 0\), we need only compute in degrees 1 and 2. By using the same standard generators for \(\tilde{H}_1(S^1)\) as part (a) and the corresponding induced generators on \(\tilde{H}_1(T^2)\), we have that \(\iota_*: \tilde{H}_1(S^1) \to \tilde{H}_1(S^1) \oplus \tilde{H}_1(T^2)\) is identified with \(\mathbb{Z} \to \mathbb{Z} \oplus (\mathbb{Z} \oplus \mathbb{Z}), 1 \mapsto (2,(0,1))\) where the multiplication by 2 is induced by including to the boundary of the mobius strip being the winding twice around map under the homotopy identifying \(M \simeq S^1\). This shows that \(\iota_*\) is injective so we find that \(\delta: \tilde{H}_2(X) \to \tilde{H}_1(S^1)\) is the zero map, and thus \(\tilde{H}_2(X) \cong \tilde{H}_2(T^2) \cong \mathbb{Z}\) by exactness. It remains to compute in degree 1. Now we use that \(j_*: \tilde{H}_1(S^1) \oplus \tilde{H}_1(T^2) \to \tilde{H}_1(X)\) satisfies \(\ker j_* = \operatorname{Im} \iota_* = \mathbb{Z}(2,0,1)\), it follows that \(j_* \mathbb{Z}^3 = \mathbb{Z}/2 \mathbb{Z} \oplus \mathbb{Z}\), but since \(\tilde{H}_0(S^1) = 0\), we have \(\delta: \tilde{H}_1(X) \to \tilde{H}_0(S^1)\) satisfies \(\ker \delta = \tilde{H}_1(X)\), so that by exactness \(\tilde{H}_1(X) = \operatorname{Im}j_* \cong \mathbb{Z}^3/\mathbb{Z}(2,0,1)\), since we can take the following as a basis for \(\mathbb{Z}^3\): \(\set{(2,0,1),(1,0,0),(0,1,0)}\), we find that written in this basis \(\tilde{H}_1(X) \cong \mathbb{Z}^3/\mathbb{Z} \oplus 0 \oplus 0 \cong \mathbb{Z}^2\). Everything taken together gives us
        \begin{align*}
            \tilde{H}_k(X) = \begin{cases}
                \mathbb{Z} & k = 2 \\ \mathbb{Z} \oplus \mathbb{Z} & k = 1 \\ 0 & \text{else}
            \end{cases}
        \end{align*}
        \qed
    \end{pb}
    \begin{pb}
        \textbf{(a)} We use homework 5 problem 1, that a map of short exact sequences of chain complexes descends on homology to a map of the long exact sequences given by the snake lemma. This is proven in homework 5. We need a new naturality property, i.e. naturality of the good pair, this means that suppose we have \((X,A)\) a good pair and \(f: (X,A) \to (X,A)\) satisfying that the following commutes, for \(f'\) the induces map on the quotient
        \begin{equation*}
            \begin{tikzcd}
                A \arrow[d,"f"] \arrow[r,"\iota"] &X \arrow[d,"f"] \arrow[r,"\pi"] &X/A \arrow[d,"f'"] \\
                A \arrow[r,"\iota"] &X \arrow[r,"\pi"] &X/A
            \end{tikzcd}
        \end{equation*}
        Then given the following SES of short exact sequences, where once again \(\hat{f}\) is the induces map
        \begin{equation*}
            \begin{tikzcd}
                0 \arrow[r,""] &C_*(A) \arrow[d,"f_*"] \arrow[r,"\iota_*"] &C_*(X) \arrow[d,"f_*"] \arrow[r,"q"] &C_*(X,A) \arrow[d,"\widehat{f}"] \arrow[r,""] &0 \\
                0 \arrow[r,""] &C_*(A) \arrow[r,"\iota_*"] &C_*(X) \arrow[r,"q"] &C_*(X,A) \arrow[r,""] &0
            \end{tikzcd}
        \end{equation*}
        The induced maps on homology from the functoriality of snake lemma satisfy \([\hat{f}] = [f']\) when \(\tilde{H}_*(X,A)\) is identified with \(\tilde{H}_*(X/A)\).
        To prove this, it suffices to show that \(\pi: (X,A) \to (X/A,A/A)\) induces the isomorphism on homology, since then we get \(\pi_* [q] = \pi_*\), and thus \([\hat{f}] = \pi_*[q][f] = \pi_*[f] = [f']\). To see that \(\pi: (X,A) \to (X/A,A/A)\) induces the isomorphism on homology

        \textbf{(b)} Let \(f: X \to X\), then we can define \(Cf: CX \to CX\) as \(\pi_{CX} \circ (f \times 1_I)\), and \(\Sigma f: \Sigma X \to \Sigma X\) as \(\pi_{\Sigma X} \circ (f \times 1)\) where \(\pi_{CX},\pi_{\Sigma X}\) are the quoients \(X \times I \to CX \tand \Sigma X\) respectively. This gives us the following commutative diagram:
        \begin{equation*}
            \begin{tikzcd}
                X \arrow[d,"f"] \arrow[r,"\iota"] &CX \arrow[d,"Cf"] \arrow[r,"\pi"] &\Sigma X \arrow[d,"\Sigma f"] \\
                X \arrow[r,"\iota"] &CX \arrow[r,"\pi"] &\Sigma X
            \end{tikzcd}
        \end{equation*}
        Moreover, we get the following SES of chain complexes
        \begin{equation*}
            \begin{tikzcd}
                0 \arrow[r,""] &C_*(X) \arrow[d,"f_*"] \arrow[r,"\iota_*"] &C_*(CX) \arrow[d,"Cf_*"] \arrow[r,"q"] &C_*(CX,X) \arrow[d,"\widehat{Cf}"] \arrow[r,""] &0 \\
                0 \arrow[r,""] &C_*(X) \arrow[r,"\iota_*"] &C_*(CX) \arrow[r,"q"] &C_*(CX,X) \arrow[r,""] &0
            \end{tikzcd}
        \end{equation*}
        Using the solution to (a) gives us
        \begin{equation*}
            \begin{tikzcd}
                H_{n+1}(\Sigma X)
            \end{tikzcd}
        \end{equation*}
    \end{pb}

    \begin{pb}
        \textbf{(a)} Assume for contradiction \(f(x) \neq \pm x\) for any \(x\), embed \(S^n\) into \(\mathbb{R}^{n+1}\), then define \(\pi(x,y)\) to be the projection of \(y\) onto \(T_xS^n = (\mathbb{R} x)^{\perp} \subset \mathbb{R}^{n+1}\), then \(\pi(x,y) : \mathbb{R}^{n+1} \times \mathbb{R}^{n+1} \to \mathbb{R}^{n+1}\) is continuous since it can be written as \((x,y) \mapsto y - \gen{x,y}x\), this implies the following is a composition of continuous functions and thus continuous
        \begin{align*}
            x \mapsto \pi(x,f(x) - x)
        \end{align*}
        This is nonvanishing since \(f(x) \neq \pm x\), so this map gives a nonvanishing continuous vectorfield on \(S^n\), this contradicts the hairy ball theorem since \(n\) is even. \qed

        \textbf{(b)} Suppose \(f: \mathbb{RP}^n \to \mathbb{RP}^n\) for \(n\) even, then \(F: S^n \to \mathbb{RP}^n\) via \(x \mapsto f\circ\rho(x)\) where \(\rho\) is the projection is continuous by definition of the quotient topology. Now since \(S^n\) is a covering space of \(\mathbb{RP}^n\), the map \(F\) has a lift to \(\widehat{F}: S^n \to S^n\) satisfying \(\rho\circ \widehat{F} = F\), the construction of each of these maps gives commutativity of
        \begin{equation*}
            \begin{tikzcd}
                S^n \arrow[d,"\rho"] \arrow[r,"\widehat{F}"] & S^n \arrow[d,"\rho"] \\
                \mathbb{RP}^n \arrow[r,"f"] &\mathbb{RP}^n
            \end{tikzcd}
        \end{equation*}
        Now, by part (a) \(\widehat{F}(x) = \pm x\) for some \(x \in S^n\), so that \(f([x]) = [\widehat{F}(x)] = [x]\). \qed
    \end{pb}
\end{document}